%% kapa: exam=1 round=23 year=2012 period=2 form=A subject=accounting subject_ko=회계학 seq=23 number=23
%% source: Q-Net 감정평가사 기출문제 (cst00309) artlSeq=5052580 file=raw/1008/5052580/2012년 제23회 감정평가사 제1차시험 기출문제(2교시 A형).pdf
%% conversion: kordoc 4.24.51 -> tools/gen1.py (auto); review=proofread-sonnet; flags=none
\begin{question}{23}
20×1년 1월 1일에 (주)한국은 20×1년 12월 31일부터 매년 말 ￦100,000씩 3년간 수취하는 조건으로 상품을 할부판매하였다. 이 상품의 현금판매가격은 ￦257,710으로 취득 당시의 유효이자율 8\%를 반영하여 결정된 것이다. 유효이자율법을 적용하여 회계처리하는 경우, 20×1년 12월 31일 판매대금 ￦100,000을 회수할 때 인식하여야 하는 이자수익은? (단, 계산금액은 소수점 첫째자리에서 반올림한다.)
\begin{choices}
  \item ￦8,000
  \item ￦16,000
  \item ￦20,617
  \item ￦22,266
  \item ￦24,000
\end{choices}
\end{question}
\answer{3}
